\pdfinfoomitdate=1
\pdftrailerid{}
\pdfsuppressptexinfo=15
\documentclass[11pt,letterpaper]{article}
\usepackage[margin=0.94in]{geometry}
\usepackage[T1]{fontenc}
\usepackage{lmodern,amsmath,amssymb,amsthm,mathtools,booktabs,microtype}
\usepackage[colorlinks=true,linkcolor=black,citecolor=black,urlcolor=blue]{hyperref}
\usepackage{enumitem,fancyhdr,xcolor}
\setlist{itemsep=3pt,topsep=5pt}
\setlength{\parskip}{5pt}\setlength{\parindent}{0pt}
\setlength{\emergencystretch}{2em}
\allowdisplaybreaks[2]
\newtheorem{theorem}{Theorem}[section]
\newtheorem{lemma}[theorem]{Lemma}
\newtheorem{proposition}[theorem]{Proposition}
\newtheorem{corollary}[theorem]{Corollary}
\theoremstyle{remark}\newtheorem{remark}[theorem]{Remark}
\newcommand{\Prob}{\mathbb P}\newcommand{\E}{\mathbb E}
\newcommand{\U}{\mathcal U}\newcommand{\dd}{\,\mathrm d}
\newcommand{\rf}[2]{#1^{\overline{#2}}}
\newcommand{\co}[2]{[#1^{#2}]}
\DeclareMathOperator{\Aut}{Aut}
\pagestyle{fancy}\fancyhf{}\fancyhead[L]{\small Report 131\quad Leafless multigraphs}\fancyhead[R]{\small 2 October 2026}\fancyfoot[C]{\thepage}
\setlength{\headheight}{14pt}
\hypersetup{pdftitle={Report 131: Leafless loopless multigraphs by edges},pdfauthor={}}
\begin{document}
\thispagestyle{empty}
\begin{center}
{\small\sffamily REPORT 131}\par\vspace{12pt}
{\LARGE\bfseries Leafless loopless multigraphs\par by number of edges}\par\vspace{10pt}
{\large Refined logarithmic asymptotics, connectedness,\par and inversion for OEIS A307316 and A307317}\par\vspace{12pt}
{2 October 2026}
\end{center}
\vspace{8pt}
\begin{abstract}
Let $A_m$ and $C_m$ count isomorphism classes of loopless multigraphs with $m$ edges and minimum degree at least two, with $C_m$ restricted to connected graphs. Both vertices and edges are unlabelled. With $w=W(2m)$ and
\[
 F(m)=m\left(w-\log w-1+\frac2w\right),\qquad
 H_m=\frac{\exp\{F(m)+w^2/4-3w/2-1\}}
 {\sqrt{2\pi m(w+1)}},
\]
we prove the polynomial sandwich
\[
 (1-o(1))H_m\le C_m\le A_m\le(1+o(1))e^wH_m.
\]
It yields $\log a_m=F(m)+w^2/4+O(\log m)$ for either sequence, all fixed inverse-logarithmic orders, and an implicit first-threshold inverse with additive error $O(1)$. A separate component convolution proves $C_m/A_m=1-O(w^3/m)$ directly for isomorphism classes. An elementary star-conditioning argument also gives a relative minimum-degree probability in the vertex-labelled, exact-total composition model. These statements do not establish $A_m\sim H_m$.

The only non-elementary enumeration input is Wright's classical simple-graph theorem. Its statement was checked in reproduced primary text and corroborated by an official contemporary abstract; the original 1972 scan and full 1974 proof were not obtained. No historical priority claim is made. Exact finite checks and a reproducible source package supplement the proofs; they do not certify the asymptotic theorem.
\end{abstract}
\vfill
\noindent\textbf{Reading guide.} Sections 1--7 prove the counting results and inversion. Section 8 is a separate labelled-model refinement. Section 9 explains exact checks, remaining problems, and interpretation. All logarithms are natural.
\newpage
\tableofcontents
\newpage

\section{Objects, conventions, and the main result}
A loopless multigraph on a finite vertex set has one nonnegative integer multiplicity for every unordered pair of distinct vertices. Its edge count and degrees count multiplicities. We count only vertices incident to edges, and exclude degree-one vertices. Thus every vertex of a nonempty graph has degree at least two. This convention is essential: allowing arbitrarily many isolated vertices would make an edge-indexed unlabelled count infinite.

Two such graphs are the same when a vertex bijection preserves every multiplicity. Parallel edges are not individually labelled. Let $A_m$ count all resulting classes and $C_m$ the connected classes. These are OEIS A307316 and A307317 \cite{oeisA,oeisC}. For $m>0$ the handshake identity gives at most $m$ vertices. There is no one-edge graph, and there is exactly one two-edge graph: a pair of vertices with multiplicity two. We adopt the entries' exceptional $A_0=C_0=1$ convention; $C_0$ will never enter a component product.

Write $W$ for the positive real Lambert function, $W(x)e^{W(x)}=x$, and, for real $m>0$, set
\begin{align}
 w&=W(2m),& v&=e^w=\frac{2m}{w},\label{eq:wv}\\
 F(m)&=m\left(w-\log w-1+\frac2w\right),&
 J(m)&=F(m)+\frac{w^2}{4},\label{eq:FJ}\\
 H_m&=\frac{\exp\{J(m)-3w/2-1\}}{\sqrt{2\pi m(w+1)}}.\label{eq:H}
\end{align}
All asymptotics below are as the integer $m$ tends to infinity, unless another limit is specified. An equality with an $O$ term for $a_m\in\{A_m,C_m\}$ means two separate estimates.

\begin{theorem}[Counting, ratio, and inverse]\label{thm:main}
Using Wright's simple-graph theorem stated in Section~\ref{sec:wright},
\begin{align}
 (1-o(1))H_m&\le C_m\le A_m\le(1+o(1))e^wH_m,\label{eq:sandwich}\\
 \log a_m&=J(m)+O(\log m),\qquad a_m\in\{A_m,C_m\},\label{eq:log}\\
 \frac{C_m}{A_m}&=1-O\left(\frac{w^3}{m}\right).\label{eq:ratio}
\end{align}
For $X$ sufficiently large, define the first threshold
\[
 \nu_a(X)=\min\{m\ge0:a_m\ge X\}.
\]
If $m_0$ is the unique sufficiently large real solution of $J(m_0)=\log X$, then
\begin{equation}
 \nu_a(X)=m_0+O(1).\label{eq:inverse}
\end{equation}
No monotonicity assumption on either sequence is used.
\end{theorem}

The multiplicative gap in \eqref{eq:sandwich} is $e^w=2m/w$. Therefore the theorem is not a relative equivalent for either count. The ratio conclusion is stronger than equality of logarithmic expansions and requires its own proof. Section~\ref{sec:series} gives an effective extraction procedure for every fixed inverse-logarithmic order, with explicit first terms.

\section{Classical input and the simple-graph saddle}\label{sec:wright}
Let $S_m$ count unlabelled simple graphs with $m$ edges and no isolated vertices. Padding with isolates to exactly $2m$ vertices is an isomorphism-class bijection, so $S_m=T(2m,m)$ in Wright's notation.

\subsection{The precise classical theorem used}
The stable-saddle specialization of Wright's 1972 Theorem 5 \cite{wright72} is
\begin{equation}
 S_m=(1+o(1))e^{-1}\sqrt{\frac{2\pi V}{1+\log V}}
       \frac{\binom{V(V-1)/2}{m}}{V!},
 \qquad V=\lfloor v\rfloor.\label{eq:wright}
\end{equation}
Indeed $V$ is the greatest integer satisfying $V\log V\le2m$. Wright's theorem applies when his parameter $\mu=2q/n-\log n$ is below $-B n^{-1/2}\log n$ for a fixed $B>0$. At $n=2m,q=m$, $\mu=1-\log(2m)$ satisfies this eventually. His standing restriction $q\le\binom n2/2$ also holds eventually. The published error is a relative $1+O(m^{-c})$ for some $c>0$; only $1+o(1)$ is needed here.

\textbf{Source-verification boundary.} The statement and definitions were checked in an online reproduction of the primary article \cite{wrighttext}; an official January 1972 AMS Notices abstract independently corroborates the stable saddle \cite{wrightabstract}. The original AMS scan could not be retrieved, and the full 1974 successor proof \cite{wright74} was not inspected. The reproduced text has OCR defects, so the official abstract is corroboration rather than a substitute for the precise statement. The proof below explicitly takes \eqref{eq:wright} as a classical input. A publication-oriented bibliography check should obtain an original scan or the full successor paper. Nothing here certifies novelty.

\subsection{Refining the simple-graph equivalent}
Define
\begin{equation}
 Q_m=\frac{\exp\{F(m)-w^2/4-w/2-1\}}
 {\sqrt{2\pi m(w+1)}}.\label{eq:Q}
\end{equation}
We claim that $S_m\sim Q_m$. To see all constants, initially use a real vertex parameter $n$ near $v$, and write $N=n(n-1)/2$. The falling-factorial expansion is
\begin{equation}
 \log\binom Nm=m\log N-\log\Gamma(m+1)
 -\frac{m(m-1)}{2N}+O(m^3/N^2).\label{eq:falling}
\end{equation}
This follows by summing $\log(1-j/N)=-j/N+O(j^2/N^2)$ for $0\le j<m$. Here $m/N\to0$, so the expansion is uniform in a small relative neighborhood of $v$.

At $n=v$, $m/v=w/2$ and
\[
 m\log N=m\log(v^2/2)-w/2+o(1),\qquad
 \frac{m(m-1)}{2N}=\frac{w^2}{4}+o(1).
\]
The error $m^3/v^4=O(w^3/v)$ tends to zero. Stirling's formula gives
\begin{equation}
 \log\frac{\binom{v(v-1)/2}{m}}{\Gamma(v+1)}
 =F(m)-\frac{w^2}{4}-\frac w2
 -\frac12\log(2\pi m)-\frac12\log(2\pi v)+o(1).\label{eq:simpleweight}
\end{equation}
For verification of the main exponent, define
\begin{equation}
 \phi_m(n)=m\log\frac{n^2}{2m}+m-n\log n+n.\label{eq:phi}
\end{equation}
Then
\[
 \phi_m(v)=F(m),\quad \phi_m'(v)=0,\quad
 \phi_m''(v)=-\frac{w+1}{v}.
\]
Rounding $v$ to $V$ changes the principal term by $O(w/v)$; the finite-density corrections change by $O(w^2/v)$ and the logarithmic factors by $O(1/v)$. These are all $o(1)$. Combining \eqref{eq:simpleweight} with the prefactor in \eqref{eq:wright} proves
\begin{equation}
 S_m\sim Q_m,\qquad \log S_m=F(m)+O(\log^2m).\label{eq:S}
\end{equation}
The exact derivative identities needed later are
\begin{equation}
 w'(m)=\frac{w}{m(1+w)},\quad
 F'(m)=w-\log w,\quad
 F''(m)=\frac{w-1}{m(1+w)}.\label{eq:Fderiv}
\end{equation}
In particular, $F$ is eventually increasing and convex.

\section{Uniform compositions and a rare-event lower bound}\label{sec:compositions}
For an integer $n\ge2$, let $N=\binom n2$. The model $\U(n,m)$ is uniform on all weak compositions of $m$ into the $N$ multiplicity coordinates. It has exactly
\begin{equation}
 \binom{N+m-1}{m}\label{eq:weakcount}
\end{equation}
vertex-labelled graphs. This exact count is classical, appearing, for example, in Gilbert \cite{gilbert}. Write $B=\{\delta\ge2\}$, $D=\{\text{disconnected}\}$, and $\lambda=2m/n$.

\subsection{Two elementary tools}
\begin{lemma}[Association for independent coordinates]\label{lem:harris}
For increasing events $E,F$ in a product of independent ordered random variables, $\Prob(E\cap F)\ge\Prob(E)\Prob(F)$. The same statement holds for bounded increasing functions, with covariance in place of the event inequality.
\end{lemma}
\begin{proof}
For one coordinate $X$ and an independent copy $X'$, the covariance of increasing $f,g$ equals
\[
 \frac12\E[(f(X)-f(X'))(g(X)-g(X'))]\ge0.
\]
Induct on the number of coordinates: conditional covariance in the last coordinate is nonnegative, and the two conditional expectations are increasing functions of the remaining independent coordinates. The covariance decomposition completes the induction. Indicators give the event statement. This is the product-measure association inequality usually called Harris's inequality.
\end{proof}

\begin{lemma}[Monotone composition coupling]\label{lem:couple}
If a composition $a$ is uniform at total $k$, add one to coordinate $i$ with probability $(a_i+1)/(N+k)$. The new composition is uniform at total $k+1$. Thus the probability of any increasing property is nondecreasing in the total.
\end{lemma}
\begin{proof}
A target vector $b$ of total $k+1$ receives mass
\[
 \frac1{\binom{N+k-1}{k}}\frac{\sum_i b_i}{N+k}
 =\frac1{\binom{N+k}{k+1}}.
\]
Each transition increases one coordinate and decreases none.
\end{proof}

\subsection{Geometric conditioning and quantitative lower bound}
Assign independent geometric multiplicities
\[
 \Prob(X_e=j)=(1-\rho)\rho^j,\qquad j\ge0.
\]
Every vector with total $k$ has equal probability $(1-\rho)^N\rho^k$, so conditioning on the total gives exactly $\U(n,k)$.

\begin{lemma}\label{lem:rare-lower}
If $\lambda=\log n+o(1)$, then
\begin{equation}
 \Prob_{\U(n,m)}(B)\ge(1-o(1))
 \exp\{-n(1+\lambda)e^{-\lambda}\}.\label{eq:lowerB}
\end{equation}
The estimate is uniform on every deterministic envelope
$|\lambda-\log n|\le\eta_n$ with $\eta_n\to0$.
\end{lemma}
\begin{proof}
Put $\ell=\log n$ and
\[
 \mu=m-3\sqrt{m\ell},\quad \rho=\frac\mu{N+\mu},\quad
 \lambda'=\frac{2\mu}{n}.
\]
These parameters are positive and valid eventually. The total $T=\sum_eX_e$ has mean $\mu$. The probability that one vertex has degree at most one is exactly
\[
 q=(1-\rho)^{n-1}\{1+(n-1)\rho\}.
\]
Lemma~\ref{lem:harris}, iterated over the increasing good-degree events, gives $\Prob(B)\ge(1-q)^n$ in the product model. Uniformly for $|\lambda-\ell|\le1$,
\begin{align*}
 \lambda'-\lambda&=O(\ell/\sqrt n),&
 (n-1)\rho&=\lambda'+O(\ell^2/n),\\
 (n-1)\log(1-\rho)&=-\lambda'+O(\ell^2/n),&
 nq&=n(1+\lambda')e^{-\lambda'}+O(\ell^3/n).
\end{align*}
The derivative of $n(1+x)e^{-x}$ is $-nxe^{-x}=O(\ell)$ in this range. Also $nq^2=O(\ell^2/n)$. Consequently
\begin{equation}
 (1-q)^n=\exp\{-n(1+\lambda)e^{-\lambda}+O(\ell^2/\sqrt n)\}.\label{eq:harrisquant}
\end{equation}

The centered cumulant generating function of $T$ is
\begin{equation}
 K_\mu(t)=N\log\frac{1-\rho}{1-\rho e^t}-\mu t
 =\frac{\mu(1+\mu/N)}2t^2+O(\mu|t|^3),\label{eq:mgf}
\end{equation}
uniformly for the present parameters and small $|t|$. This follows by differentiating three times: the second derivative at zero is $\mu(1+\mu/N)$, and the third derivative is $O(\mu)$ on a fixed small neighborhood when $\rho=o(1)$.

Let $d=m-\mu=3\sqrt{m\ell}$ and choose $t=d/\mu=O(n^{-1/2})$. The moment-generating function exists there, and Chernoff's inequality yields
\begin{align}
 \log\Prob(T>m)
 &\le-\frac{d^2}{2\mu}
 +O\left(\frac{d^2}{N}+\frac{d^3}{\mu^2}\right)\notag\\
 &=-\frac92\ell+O(\ell/\sqrt n+\ell^2/n).\label{eq:tail}
\end{align}
This is an upper bound, not a tail equivalent.

Let $f_k=\Prob_{\U(n,k)}(B)$. Lemma~\ref{lem:couple} implies
\[
 \Prob(B)=\E f_T\le f_m+\Prob(T>m).
\]
Under the lemma's hypothesis the exponential in \eqref{eq:lowerB} is $n^{-1+o(1)}$, while the upper bound \eqref{eq:tail} is at most $n^{-9/2+o(1)}$. Subtraction is therefore negligible on the required relative scale. All errors just used are uniform in a bounded neighborhood of $\ell$; the rare-event comparison is uniform on any shrinking envelope. This proves the claim.
\end{proof}

An additive approximation to a probability tending to zero would not suffice here. The error in \eqref{eq:harrisquant} is additive in its logarithm and tends to zero; that is why this argument preserves a relative lower bound. Section~\ref{sec:star} supplies the matching upper bound by a separate argument.

\section{Small-component bounds in the labelled model}\label{sec:cuts}
\begin{lemma}[Connectedness after the degree restriction]\label{lem:cuts}
If $\lambda=\log n+o(1)$, then
\begin{equation}
 \Prob_{\U(n,m)}(D\cap B)=O(\log^2n/n^2),\qquad
 \Prob_{\U(n,m)}(D\mid B)\le n^{-1+o(1)}.\label{eq:cuts}
\end{equation}
If additionally $(\lambda-\log n)\log n=o(1)$, the latter bound improves to $O(\log^2n/n)$.
\end{lemma}
\begin{proof}
Fix an $s$-element vertex set, where $2\le s\le n/2$. Let
\[
 c=s(n-s),\quad K=\binom s2,\quad M=N-c,\quad r=\frac m{N+m}.
\]
The exact probability that its cross-edge coordinates are all zero is
\begin{equation}
 \frac{\binom{M+m-1}{m}}{\binom{N+m-1}{m}}
 =\prod_{i=0}^{m-1}\left(1-\frac c{N+i}\right)
 \le e^{-rc}.\label{eq:cutexact}
\end{equation}
Conditionally on this empty cut, the remaining $M$ coordinates form a uniform weak composition. Let $X$ be the total in the $K$ internal coordinates. Its binomial factorial moment is
\begin{equation}
 \E\binom Xs=\binom ms\frac{\rf K s}{\rf M s}
 \le\left(\frac{em(K+s)}{sM}\right)^s.\label{eq:moment}
\end{equation}
Here $a^{\overline{s}}=a(a+1)\cdots(a+s-1)$. One derivation uses
\[
 \sum_{x\ge0}\binom xs\binom{K+x-1}{x}z^x
 =\binom{K+s-1}{s}z^s(1-z)^{-K-s}.
\]
Multiply by $(1-z)^{-(M-K)}$, extract the coefficient of $z^m$, and divide by \eqref{eq:weakcount} with $N=M$. The inequality in \eqref{eq:moment} follows from $\binom ms\le(em/s)^s$, $\rf K s\le(K+s)^s$, and $\rf M s\ge M^s$. Since $\binom Xs\ge1$ when $X\ge s$, the same bound controls that probability.

A disconnected graph in $B$ has a component of size $2\le s\le n/2$. Its cut is empty and, by the handshake identity, its internal edge count satisfies $X\ge s$. Multiplicity, not merely the number of occupied pairs, is counted. In particular this includes a two-vertex component joined by parallel edges. Union-bound \eqref{eq:cutexact} times the conditional bound \eqref{eq:moment} over all $\binom ns$ subsets.

For $2\le s\le n/\ell^2$, where $\ell=\log n$, we have $M\ge n^2/3$, $K+s\le s^2$, and $m=O(n\ell)$ eventually. Moreover
\[
 rn=\frac{\lambda n}{n-1+\lambda}=\ell+o(1),\qquad
 r(n-s)\ge\ell-1
\]
eventually. The contribution for each $s$ is at most
\[
 \left(\frac{en}{s}\right)^s
 \left(\frac en\right)^s
 \left(\frac{Cs\ell}{n}\right)^s
 \le\left(\frac{C'\ell}{n}\right)^s.
\]
Summing from $s=2$ gives $O(\ell^2/n^2)$.

For $n/\ell^2<s\le n/2$, omit the internal-edge requirement. Eventually $r(n-s)\ge0.49\ell$ and $\log(en/s)\le1+2\log\ell$, so the summand is at most $\exp(-0.4s\ell)$. Its sum is at most $n\exp(-0.4n/\ell)$, after harmless adjustment of integer endpoints. This is negligible compared with $\ell^2/n^2$.

Dividing the established bound by Lemma~\ref{lem:rare-lower} proves the conditional estimate. Under the stronger hypothesis,
\[
 n(1+\lambda)e^{-\lambda}=1+\log n+o(1),
\]
so the lower bound for $\Prob(B)$ is $(e^{-1}+o(1))/n$. This proves the stated improvement. The constants in the first bound can be uniform over $|\lambda-\log n|\le1$, replacing $\ell-1$ above by $\ell-2$.
\end{proof}

This lemma concerns a vertex-labelled uniform-composition measure. It does not by itself prove connectedness under the uniform measure on unlabelled classes. Section~\ref{sec:ratio} gives a separate, automorphism-safe argument for that conclusion.

\section{The polynomial sandwich}
\subsection{An upper bound through simple supports}\label{sec:support}
Let $M_m$ count loopless unlabelled multigraphs with $m$ edges and no isolated vertices, without excluding leaves. Collapsing each parallel family yields a simple support with $k$ edges. For each support class fix one representative and an arbitrary ordering of its edges. Positive multiplicities summing to $m$ give $\binom{m-1}{k-1}$ possibilities, mapping surjectively onto the multigraph classes with that support. Support automorphisms may identify assignments; they cannot invalidate the upper bound. Therefore
\begin{equation}
 A_m\le M_m\le\sum_{j=0}^{m-1}S_{m-j}\binom{m-1}{j}.\label{eq:support}
\end{equation}
No assumption about asymmetry or a factorial division is used.

For $0\le j\le w^3$, \eqref{eq:S} and Taylor expansion give uniformly
\begin{equation}
 \frac{S_{m-j}}{S_m}=e^{-jF'(m)+o(1)},\qquad
 \binom{m-1}{j}=\frac{m^j}{j!}(1+o(1)).\label{eq:central-support}
\end{equation}
Indeed $F''=O(1/m)$, so the quadratic remainder is $O(w^6/m)=o(1)$. The derivative of $\log Q_m-F(m)$ is $O(w/m)$ and contributes $O(w^4/m)=o(1)$. The binomial relative logarithmic error is $O(j^2/m)=o(1)$. The error $S_k/Q_k-1$ is automatically uniform whenever the least index $k$ tends to infinity. Since
\[
 m e^{-F'(m)}=mwe^{-w}=\frac{w^2}{2},
\]
the central sum in \eqref{eq:support} is
\begin{equation}
 (1+o(1))S_m e^{w^2/2}.\label{eq:supportmain}
\end{equation}
The exponential-series tail beyond $w^3$ is relatively negligible because $w^3/(w^2/2)\to\infty$.

For $w^3<j\le m/2$, convexity and the coarse estimate in \eqref{eq:S} bound the summand by
\[
 e^{F(m)+O(w^2)}\frac{a^j}{j!},\qquad
 a=m e^{-F'(m/2)}=W(m)^2\asymp w^2.
\]
For an integer $q$ within one of $w^3$, the factorial bound and consecutive-term ratio give
\[
 \sum_{j\ge q}\frac{a^j}{j!}\le2(ea/q)^q
 =\exp\{-\Omega(w^3\log w)\}.
\]
This absorbs $e^{O(w^2)}$ and is negligible relative to \eqref{eq:supportmain}. Finally, the range $m-j<m/2$ contributes at most
\[
 \exp\{F(m/2)+m\log2+O(w^2)\}.
\]
Here finitely many small support sizes are absorbed separately; the sum of the binomial weights is at most $2^{m-1}$. Since $F(m)-F(m/2)=\Theta(m\log m)$, this range is negligible too. We obtain
\begin{equation}
 A_m\le(1+o(1))Q_m e^{w^2/2}
 =(1+o(1))e^wH_m.\label{eq:upperH}
\end{equation}
In particular, $\log A_m\le F(m)+O(\log^2m)$, a coarse consequence used in the ratio proof.

\subsection{A central Gaussian lower bound}
For an integer $n$ define
\begin{equation}
 T_{n,m}=\frac1{n!}\binom{\binom n2+m-1}{m},\quad
 \sigma^2=\frac v{w+1},\quad
 I_m=\{n\in\mathbb Z:|n-v|\le w\sigma\}.\label{eq:T}
\end{equation}
A class $G$ with $n$ vertices contributes $1/|\Aut(G)|$ to labelled count divided by $n!$. Consequently $T_{n,m}$ is a symmetry-weighted count, not an unlabelled count; restricted to any property it is a valid lower bound for the number of classes having that property.

Write $n=v+h$ and $x=h/v$. On $I_m$, $x=O(\sqrt{w/v})$, and
\begin{equation}
 \lambda-\log n=\frac{w}{1+x}-w-\log(1+x)
 =-(w+1)x+O(wx^2)=O(w^{3/2}/\sqrt v).\label{eq:windowlambda}
\end{equation}
In particular $(\lambda-\log n)\log n=o(1)$ uniformly. More precisely, for $a(n)=n(1+2m/n)e^{-2m/n}$,
\[
 a(v)=1+w,\qquad
 a'(n)=e^{-\lambda}(1+\lambda+\lambda^2)=O(w^2/v),
\]
so $a(n)=1+w+O(w^{5/2}/\sqrt v)=1+w+o(1)$ throughout the window. Lemmas~\ref{lem:rare-lower} and \ref{lem:cuts} imply
\begin{equation}
 \Prob_{\U(n,m)}(B\cap D^c)
 \ge(1-o(1))e^{-1-w}\quad(n\in I_m),\label{eq:jointwindow}
\end{equation}
uniformly: the subtracted disconnected probability is $O(w^2/v^2)=o(e^{-w})$.

The rising-factorial analogue of \eqref{eq:falling}, followed by Stirling, gives
\begin{equation}
 \log T_{n,m}=\phi_m(n)+\frac{\lambda^2}{4}-\frac\lambda2
 -\frac12\log(2\pi m)-\frac12\log(2\pi n)+o(1)\label{eq:rise}
\end{equation}
uniformly on $I_m$. The sign of the $\lambda^2/4$ term is positive: repeated pairs are permitted. Its remainder is $O(m^3/n^4)=o(1)$.

The third derivative of \eqref{eq:phi} is $4m/n^3+1/n^2=O(w/v^2)$. Since $|h|\le w\sigma=O(\sqrt{wv})$, the cubic remainder is $O(w^{5/2}/\sqrt v)=o(1)$. Changes in $\lambda^2/4-\lambda/2$ and $\log n$ are also $o(1)$. Thus
\begin{equation}
 T_{n,m}=(1+o(1))
 \frac{e^{F(m)+w^2/4-w/2}}{\sqrt{2\pi m}\sqrt{2\pi v}}
 \exp\left\{-\frac{(n-v)^2}{2\sigma^2}\right\}\label{eq:gauss}
\end{equation}
uniformly on $I_m$. As $\sigma\to\infty$ and $w\to\infty$, a shifted-lattice Riemann sum gives
\[
 \sum_{n\in I_m}\exp\left\{-\frac{(n-v)^2}{2\sigma^2}\right\}
 =(1+o(1))\sqrt{2\pi}\sigma.
\]
For completeness, first restrict standardized deviations to a fixed compact interval and use Riemann sums of mesh $1/\sigma$; the complementary sums are bounded by the integrable Gaussian tails. Then let the compact interval increase. The fractional part of $v$ is immaterial.

Graphs with different vertex counts are different isomorphism classes. Multiplying \eqref{eq:gauss} by \eqref{eq:jointwindow}, then summing the symmetry-weighted lower bounds, proves
\[
 C_m\ge(1-o(1))e^{-1-w}
 \frac{e^{F(m)+w^2/4-w/2}}{\sqrt{2\pi m(w+1)}}
 =(1-o(1))H_m.
\]
Together with \eqref{eq:upperH} and $C_m\le A_m$, this proves \eqref{eq:sandwich}. Only a central-window lower bound was needed; no assertion about omitted global vertex-count tails was used.

\section{The connected-to-total ratio for isomorphism classes}\label{sec:ratio}
A nonempty loopless component of minimum degree at least two has at least two edges. Every disconnected graph counted by $A_m$ has a component with $2\le k\le m/2$ edges. Taking that component and its complement gives a surjection from a subset of pairs of leafless classes, so
\begin{equation}
 0\le A_m-C_m\le\sum_{k=2}^{\lfloor m/2\rfloor}A_kA_{m-k}.\label{eq:convolution}
\end{equation}
Repeated components only cause overcounting, which is harmless.

Let $g(x)=\log H_x$, with $w=W(2x)$. Direct differentiation yields
\begin{align}
 g'(x)&=F'(x)+\frac{w^3-3w^2-6w-1}{2x(w+1)^2}
       =F'(x)+O(w/x),\label{eq:gprime}\\
 g''(x)&=\frac{w-1}{x(w+1)}+O(w/x^2)>0
 \quad\hbox{eventually}.\label{eq:gsecond}
\end{align}
Set $t=e^{-g'(m/2)}$. With $u=W(m)$, $t\sim u^2/m\sim w^2/m$. The upper sandwich, uniformly for $m/2\le m-k\le m$, and convexity give
\begin{equation}
 \frac{A_{m-k}}{H_m}
 \le(1+o(1))e^w e^{g(m-k)-g(m)}
 \le(1+o(1))v t^k.\label{eq:ratio-middle}
\end{equation}

We now prove
\begin{equation}
 \sum_{k=2}^{\lfloor m/2\rfloor}A_kt^k=(1+o(1))t^2.\label{eq:small-series}
\end{equation}
The coarse bound $\log A_k\le F(k)+O(\log^2k)$ implies $A_k\le(Ck)^k$ for all $k\ge2$, after enlarging a fixed $C$ for small indices. Put $L_0=\lceil10\log m\rceil$. Since $CL_0t\to0$,
\[
 \sum_{3\le k\le L_0}A_kt^k
 \le\sum_{k\ge3}(CL_0t)^k
 =O(L_0^3t^3)=o(t^2),
\]
because $L_0^3t=O(\log^5m/m)\to0$.

For larger $k$, use
\[
 \frac{\dd}{\dd k}\frac{F(k)}k
 =\frac{1-2/W(2k)}k>0
\]
eventually. For $L_0\le k\le m/2$ this gives
\[
 \frac{F(k)}k\le\frac{F(m/2)}{m/2}
 =F'(m/2)-1+\frac2u.
\]
Hence, uniformly on this range,
\[
 \log(A_kt^k)\le k\{-1+2/u+o(1)\}+O(\log^2k)\le-k/2
\]
eventually. The uniformity follows from
$\sup_{k\ge L_0}\log^2k/k=O(\log^2L_0/L_0)=o(1)$ and \eqref{eq:gprime}. Summation gives $O(e^{-L_0/2})=O(m^{-5})=o(t^2)$. Finally $A_2=1$, proving \eqref{eq:small-series}.

Divide \eqref{eq:convolution} by $A_m\ge(1-o(1))H_m$, and use \eqref{eq:ratio-middle}--\eqref{eq:small-series}. Then
\[
 0\le\frac{A_m-C_m}{A_m}\le(1+o(1))vt^2
 =O(w^3/m).
\]
This establishes \eqref{eq:ratio} directly for unlabelled classes, without transferring the labelled probability measure or making an automorphism assumption.

\section{Logarithmic expansions and inversion}\label{sec:series}
\subsection{The extra quadratic logarithmic correction}
Using $m=we^w/2$, \eqref{eq:H} gives the exact identity
\begin{equation}
 \log H_m=J(m)-2w-\log w-1-\frac12\log\pi
            -\frac12\log(1+1/w).\label{eq:logH}
\end{equation}
The upper sandwich logarithm adds $w+o(1)$. Therefore both counting logarithms equal
\[
 \boxed{\log a_m=F(m)+\frac{W(2m)^2}{4}+O(\log m).}
\]
The extra $W^2/4$ term is larger than the remaining error. Expanding only to a fixed inverse-logarithmic order in the scale $m$ would hide this additional information.

\subsection{Every fixed forward order}
Let $L=\log(2m)$ and $\ell=\log L$. Since $w+\log w=L$, introduce the formal variable $z=L^{-1}$ and write
\[
 w=L-\ell+A(z),\qquad A(z)=\sum_{j\ge1}a_j(\ell)z^j.
\]
The defining equation is
\begin{equation}
 A+\log(1-\ell z+zA)=0.\label{eq:Aformal}
\end{equation}
It determines polynomials recursively:
\begin{equation}
 a_r(\ell)=-\co zr\log\left(1-\ell z+
 z\sum_{j<r}a_j(\ell)z^j\right),\qquad r\ge1.\label{eq:Arecur}
\end{equation}
Because $F(m)/m=2w-L-1+2/w$, its coefficients are
\begin{equation}
 \frac{F(m)}m=L-2\ell-1+\sum_{j\ge1}f_j(\ell)L^{-j},\qquad
 f_j=2a_j+2\co z{j-1}(1-\ell z+zA)^{-1}.\label{eq:Frecur}
\end{equation}
For example, $a_1=\ell$, $a_2=\ell(\ell-2)/2$, and
$a_3=\ell(2\ell^2-9\ell+6)/6$. The explicit forward terms are
\begin{align}
 \log a_m=m\bigg[&L-2\ell-1+\frac{2\ell+2}{L}
 +\frac{\ell^2}{L^2}+\frac{\ell^2(2\ell-3)}{3L^3}
 +O\left(\frac{\ell^4}{L^4}\right)\bigg].\label{eq:forward}
\end{align}
For a version retaining both scales, replace the $O$ term in brackets by the recursively continued finite expansion and its truncation error, and append $W(2m)^2/4+O(\log m)$ outside the brackets.

Here is why coefficient extraction is an asymptotic proof, rather than merely formal manipulation. Truncate \eqref{eq:Aformal} after order $K$. Taylor expansion of $\log(1+x)$, with $x=O(\ell/L)$, bounds its residual by $O_K(\ell^{K+1}/L^{K+1})$. The derivative of $w+\log w$ is $1+1/w$, bounded away from zero, so the real root differs by at most that residual. Substitution into the smooth expression for $F/m$ gives the corresponding remainder for \eqref{eq:Frecur}. The counting error $O(\log^2m)$, including the $w^2/4$ term if it is omitted, is smaller than $m/L^K$ for every fixed $K$. Thus every fixed inverse-logarithmic order is determined. This does not assert convergence of an infinite expansion or a full transseries.

In particular, the first three additive scales are
\[
 m\log m-2m\log\log m+(\log2-1)m,
\]
and $\log A_m\sim\log C_m\sim m\log m$.

\subsection{An implicit inverse with additive error \texorpdfstring{$O(1)$}{O(1)}}
By \eqref{eq:Fderiv},
\begin{equation}
 J'(m)=w-\log w+\frac{w^2}{2m(w+1)}\sim\log m.\label{eq:Jprime}
\end{equation}
It is eventually positive. Let $y=\log X$ and let $J(m_0)=y$. A scalar form convenient for computation is
\begin{equation}
 e^{w_0}(w_0^2-w_0\log w_0-w_0+2)+\frac{w_0^2}{2}=2y,
 \qquad m_0=\frac{w_0e^{w_0}}2.\label{eq:scalarJ}
\end{equation}

To prove the threshold assertion without sequence monotonicity, choose $B_0$ so that
$|\log a_m-J(m)|\le B_0\log m$ eventually. For $m_0/2\le m\le m_0-K$, eventual derivative bounds give
\[
 J(m_0)-J(m)\ge\frac12(m_0-m)\log m_0.
\]
A fixed sufficiently large $K$ makes this exceed $B_0\log m$. For all sufficiently large $m<m_0/2$, eventual increase of $J$ and
$J(m_0)-J(m_0/2)=\Theta(m_0\log m_0)$ exclude the entire range. The remaining finitely many terms cannot exceed arbitrarily large $X$. Conversely, at the integer $r=\lceil m_0+K\rceil$, the same derivative bound yields $\log a_r\ge y$ for $K$ sufficiently large. This brackets the first threshold within $O(1)$ of $m_0$ and proves \eqref{eq:inverse}.

\subsection{Explicit inverse orders and their precise error}
Let $M=F^{-1}(y)$ be the sufficiently large inverse of $F$. Write $U=\log(2y)$ and $u=\log U$. Its first terms are
\begin{align}
 M=\frac yU\bigg[1+\frac{3u+1}{U}
 +\frac{9u^2-u-2}{U^2}
 +\frac{54u^3-47u^2-26u-1}{2U^3}
 +O\left(\frac{u^4}{U^4}\right)\bigg].\label{eq:inverse-series}
\end{align}
The same fixed-order expansion applies to $m_0$ and to $\nu_a(X)$, because their difference from $M$ is $O(\log y)$, smaller than every fixed displayed expansion scale.

For effective extraction at arbitrary fixed order, put $z=1/U$, write
\[
 W(2M)=U-2u+B(z),\qquad Z=1-2uz+zB(z),\quad B(0)=0.
\]
The scalar equation $F(M)=y$ is equivalent to
\begin{equation}
 B+2\log Z+
 \log\left(1-\frac{z(u+\log Z+1)}Z+\frac{2z^2}{Z^2}\right)=0.\label{eq:Bformal}
\end{equation}
To obtain the coefficient $[z^r]B$, substitute all previously obtained coefficients in the two logarithms and negate their $z^r$ coefficient. The new coefficient enters those logarithms only at later order, so this is triangular. Then
\begin{equation}
 M=\frac yU R(z),\qquad
 R(z)=\left(Z-z(u+\log Z)-z+\frac{2z^2}Z\right)^{-1}.\label{eq:Rformal}
\end{equation}
This gives \eqref{eq:inverse-series} and every further fixed order. Taylor remainders in \eqref{eq:Bformal}, whose derivative with respect to $B$ tends to one, bound the error after any fixed truncation; substitution in \eqref{eq:Rformal} gives its stated order. Thus the procedure has analytic remainder control.

One can also retain the effect of the extra $w^2/4$ term without truncating the large-scale inverse. With $\omega=W(2M)$,
\begin{equation}
 m_0=M-\frac{\omega^2}{4(\omega-\log\omega)}+o(1),\qquad
 \nu_a(X)=M-\frac{\omega^2}{4(\omega-\log\omega)}+O(1).\label{eq:inverse-shift}
\end{equation}
Indeed $J(M)-y=\omega^2/4$, while $J'\asymp\omega$ near $M$, so $m_0-M=O(\omega)$. On this interval $F''=O(1/M)$ and the derivative of $W(2m)^2/4$ is $O(\omega/M)$. Taylor's theorem makes the residual after the displayed first-order displacement $O(\omega^2/M)=o(1)$; division by $F'(M)\asymp\omega$ is harmless. This proves \eqref{eq:inverse-shift}.

\textbf{Accuracy warning.} The additive $O(1)$ guarantee belongs to \eqref{eq:scalarJ}, or to \eqref{eq:inverse-shift} using the untruncated implicit $M$. Any fixed truncation of \eqref{eq:inverse-series} has its own much larger error, for example $O(yu^4/U^5)$ at the displayed order. It does not inherit additive $O(1)$ accuracy.

\section{An independent relative rare-event refinement}\label{sec:star}
This section proves a stronger statement for the vertex-labelled product and exact-total models. The upper probability bound developed here was not required for Theorem~\ref{thm:main}; its lower bound was already supplied by Lemma~\ref{lem:rare-lower}.

\subsection{Exact star-conditioning comparison}
\begin{theorem}\label{thm:star}
Let $n\ge3$, and give the pairs of distinct vertices independent, identically distributed nonnegative integer multiplicities. Let $B_i$ be the event that vertex $i$ has degree at most one. Put
\[
 q=\Prob(B_i),\quad \rho=\Prob(X_e>0),\quad
 q_0=\Prob(X_1+\cdots+X_{n-2}\le1).
\]
Assume $q_0<1$, and define $c=1+\rho q_0/(1-q_0)$. For every vertex $i$ and $J\subseteq[n]\setminus\{i\}$,
\begin{equation}
 q c^{-|J|}\le
 \Prob\left(B_i\,\middle|\,\bigcap_{j\in J}B_j^c\right)\le q.\label{eq:starconditional}
\end{equation}
Every conditioning event is positive. Consequently
\begin{equation}
 (1-q)^n\le\Prob(\delta\ge2)
 \le\prod_{j=0}^{n-1}(1-qc^{-j}).\label{eq:starsandwich}
\end{equation}
\end{theorem}
\begin{proof}
Nonnegativity gives $q\le q_0<1$. If $q=0$, every vertex is good almost surely and the result is immediate; do not condition on $B_i$ in this case. Hence assume $q>0$.

Delete the star at $i$ and call its variables $S$. The nonstar variables $Y$ remain independent of $S$. For $j\ne i$, let $D_j^0$ be its nonstar degree; it has $n-2$ summands, so $\Prob(D_j^0\ge2)=1-q_0$. For $K\subseteq J$ write
\[
 A_K^0=\bigcap_{j\in K}\{D_j^0\ge2\},\qquad A_\varnothing^0=\Omega.
\]
Association in the unconditional nonstar product measure gives, for $j\in K$,
\[
 \Prob(A_K^0)\ge(1-q_0)\Prob(A_{K\setminus\{j\}}^0).
\]
Thus $\Prob(A_J^0)\ge(1-q_0)^{|J|}>0$, and removing constraints in any $H\subseteq J$ gives
\begin{equation}
 \Prob(A_{J\setminus H}^0)\le\Prob(A_J^0)(1-q_0)^{-|H|}.\label{eq:remove}
\end{equation}

For a fixed star configuration $s$, let $H(s)=\{j\in J:s_{ij}>0\}$ and
\[
 A(s)=\bigcap_{j\in J}\{D_j^0+s_{ij}\ge2\}.
\]
The actual event depends on the star values, not just their support. Nonetheless nonnegativity gives
\[
 A_J^0\subseteq A(s)\subseteq A_{J\setminus H(s)}^0.
\]
If $s_{ij}=1$, dropping that endpoint's constraint is an upper bound; if $s_{ij}\ge2$, the constraint really disappears. Both cases are covered. By \eqref{eq:remove},
\[
 \Prob_Y(A(s))\le\Prob(A_J^0)(1-q_0)^{-|H(s)|}.
\]
Let $A=\bigcap_{j\in J}B_j^c$. In the unconditional star law the support indicators are independent Bernoulli variables of parameter $\rho$. Averaging gives
\begin{equation}
 \Prob(A)\le\Prob(A_J^0)
 \left(1-\rho+\frac\rho{1-q_0}\right)^{|J|}
 =\Prob(A_J^0)c^{|J|}.\label{eq:average}
\end{equation}
Also $A_J^0\subseteq A$ and $A_J^0$ is independent of $B_i$, so
$\Prob(A\mid B_i)\ge\Prob(A_J^0)$. Bayes's formula and \eqref{eq:average} now give the left inequality in \eqref{eq:starconditional}. Its right inequality follows from association because $B_i$ is decreasing and $A$ is increasing. Positivity follows from $A_J^0\subseteq A$.

Finally order the vertices. At stage $j+1$ the conditional probability of its good event lies between $1-q$ and $1-qc^{-j}$. Multiplying these positive conditional probabilities proves \eqref{eq:starsandwich}.
\end{proof}

The guard $q_0<1$ cannot be dropped from this formula: for Bernoulli edge variables at $n=3$, $q_0=1$ and $c$ is undefined. The geometric asymptotic regime below satisfies the guard.

\subsection{Logarithmic error and the geometric law}
Theorem~\ref{thm:star} implies the explicit bounds
\begin{equation}
 -nq-\frac{nq^2}{2(1-q)}
 \le\log\Prob(\delta\ge2)
 \le-nq+\frac{n(n-1)}2q\log c.\label{eq:starlog}
\end{equation}
For the lower bound use $\log(1-q)\ge-q-q^2/[2(1-q)]$. For the upper one, sum $\log(1-qc^{-j})\le-qc^{-j}$ and use $1-e^{-j\log c}\le j\log c$. Since $\log c\le\rho q_0/(1-q_0)$, when $q,q_0,\rho=O(\log n/n)$,
\begin{equation}
 \log\Prob(\delta\ge2)=-nq+O(nq^2+n^2q\rho q_0)
 =-nq+O(\log^3n/n).\label{eq:starerror}
\end{equation}

For geometric multiplicities, the mean degree is $\lambda=(n-1)\rho/(1-\rho)$ and
\[
 q=(1-\rho)^{n-1}(1+(n-1)\rho),\qquad
 q_0=(1-\rho)^{n-2}(1+(n-2)\rho).
\]
When $\lambda=\log n+o(1)$, expansion of
$\rho=\lambda/(n-1+\lambda)$ gives $q,q_0,\rho=O(\log n/n)$ and
\[
 nq=n(1+\lambda)e^{-\lambda}+O(\log^3n/n).
\]
These estimates are uniform for $|\lambda-\log n|\le1$. Therefore
\begin{equation}
 \Prob_{\mathrm{geom}}(\delta\ge2)
 =\exp\{-n(1+\lambda)e^{-\lambda}+O(\log^3n/n)\}.\label{eq:geomrelative}
\end{equation}
The logarithmic error tends to zero, yielding a relative equivalent for this rare event.

\subsection{Two-sided transfer to an exact total}
\begin{corollary}\label{cor:exactrelative}
If $2m/n=\log n+o(1)$, then
\begin{equation}
 \Prob_{\U(n,m)}(\delta\ge2)
 \sim\exp\{-n(1+2m/n)e^{-2m/n}\}.\label{eq:exactrelative}
\end{equation}
The relative error tends to zero uniformly on every shrinking envelope $|2m/n-\log n|\le\eta_n$, $\eta_n\to0$.
\end{corollary}
\begin{proof}
Put $\ell=\log n$, $d=3\sqrt{m\ell}$, $\mu_\pm=m\pm d$, and choose product geometric laws with parameters $\rho_\pm=\mu_\pm/(N+\mu_\pm)$. Let $T_\pm$ be their totals and $P_\pm$ their good-degree probabilities. For the lower-mean upper tail use $t=d/\mu_-$ in \eqref{eq:mgf}; for the upper-mean lower tail use $t=-d/\mu_+$. Both have magnitude $O(n^{-1/2})$. The two signed Chernoff calculations give
\begin{equation}
 \tau_-:=\Prob(T_->m)\le n^{-9/2}e^{o(1)},\qquad
 \tau_+:=\Prob(T_+<m)\le n^{-9/2}e^{o(1)}.\label{eq:both-tails}
\end{equation}
These are upper bounds only. The expansion is the same as \eqref{eq:tail}, with $\mu_\pm$ in place of $\mu$.

By Lemma~\ref{lem:couple}, $f_k=\Prob_{\U(n,k)}(B)$ is nondecreasing. Consequently
\[
 P_-\le f_m+\tau_-,\qquad
 P_+\ge f_m(1-\tau_+),
\]
and hence
\begin{equation}
 P_--\tau_-\le f_m\le\frac{P_+}{1-\tau_+}.\label{eq:decondition}
\end{equation}
The shifted mean degrees $\lambda_\pm=2\mu_\pm/n$ differ from $\lambda=2m/n$ by $O(\ell/\sqrt n)$. Since the derivative of $a(x)=n(1+x)e^{-x}$ is $O(\ell)$ in magnitude,
\[
 P_\pm=e^{-a(\lambda)}
 \left(1+O(\ell^2/\sqrt n+\ell^3/n)\right)
\]
by \eqref{eq:geomrelative}. Moreover $a(\lambda)=(1+o(1))\ell$, so $e^{-a(\lambda)}=n^{-1+o(1)}$. The subtracted tail in \eqref{eq:decondition} has relative size at most $n^{-7/2+o(1)}$; the upper denominator tends to one. This proves \eqref{eq:exactrelative}.

All estimates have uniform constants in a bounded neighborhood of $\ell$. Enlarging a shrinking envelope by the shifts $O(\ell/\sqrt n)$ still produces a shrinking envelope, and $a(\lambda)/\ell=1+O(\eta_n+1/\ell)$ there. This proves the stated uniformity.
\end{proof}

Combining the corollary with Lemma~\ref{lem:cuts} gives connectedness with conditional probability at least $1-n^{-1+o(1)}$ in this labelled model, and the stronger $1-O(\log^2n/n)$ bound in the narrower regime. Neither statement supplies the expected automorphism contribution required for a sharp unlabelled equivalent.

\section{Exact checks, limitations, and further questions}
\subsection{Finite checks and the component identity}
The initial values are
\[
\begin{array}{c|rrrrrrr}
 m&0&1&2&3&4&5&6\\\hline
 A_m&1&0&1&2&5&11&34\\
 C_m&1&0&1&2&4&9&26
\end{array}
\]
The reproducible package contains OEIS fixtures through $m=50$, with their public source URLs and retrieval provenance. A finite exhaustive check through six edges enumerates edge-multiplicity vectors with minimum degree two and identifies vertex-isomorphism classes canonically. This checks definitions and initial values independently of the asymptotic argument.

Unique decomposition into connected components gives the formal Euler product
\begin{equation}
 \sum_{m\ge0}A_mz^m=\prod_{k\ge1}(1-z^k)^{-C_k}.\label{eq:euler}
\end{equation}
Only $k\ge1$ occurs, so the exceptional $C_0=1$ convention causes no factor. The product is a well-defined formal series: only finitely many component sizes contribute to any coefficient. Each of the $C_k$ distinct connected types can occur with any nonnegative multiplicity, independently. Taking a formal logarithmic derivative gives
\begin{equation}
 mA_m=\sum_{j=1}^m b_jA_{m-j},\qquad
 b_j=\sum_{d\mid j}dC_d.\label{eq:eulerrec}
\end{equation}
This recurrence and its divisor inversion check every supplied fixture term.

Other exact checks evaluate the composition factorial moments, incoming masses of the monotone coupling, and rational formal-series coefficients. Explicit exception guards remain active under \texttt{python -O}. Selected intentional-corruption tests check that malformed data and altered package contents are rejected. A closed byte inventory, immutable-copy replay, normal and optimized runs, two clean byte-identical PDF builds, and deterministic ZIP extraction/repacking support reproducibility. The README specifies dependencies and scope.

\textbf{What the checks establish.} They test finite identities, transcription, implementation, and artifact integrity. They do not establish uniform asymptotic estimates, prove Wright's theorem, certify a global historical priority claim, or turn an unproved relative enumeration into a theorem.

\subsection{What remains unresolved here}
\begin{enumerate}
\item \textbf{A sharp unlabelled equivalent.} The factor-$e^w$ gap in \eqref{eq:sandwich} remains. Even the relative labelled probability \eqref{eq:exactrelative} does not prove $A_m\sim H_m$ or $C_m\sim H_m$. A sharp analysis must control global vertex-count tails and the relevant automorphism weights, or use another direct unlabelled enumeration.
\item \textbf{A leading disconnected correction.} The proved deficit is $O(w^3/m)$, obtained from a polynomial sandwich. It need not be the true order. The two-edge parallel-pair component is a natural candidate for the leading contribution, but replacing the coarse ratio control by $A_{m-2}/A_m$ asymptotics requires new work.
\item \textbf{Growing degree restrictions.} The product comparison suggests analogous questions for minimum degree at least a fixed $r>2$. No such relative theorem, nor uniformity for growing $r$, is claimed here.
\item \textbf{Full expansions.} Every fixed order in inverse powers of $\log m$ is available for the logarithm. This is distinct from relative corrections to $A_m$, exponentially small contributions, or a full transseries. The $W^2/4$ correction and implicit additive-$O(1)$ inverse are finer information, but do not close these gaps.
\item \textbf{Complete historical verification.} Wright's 1972 result and 1974 successor must be credited; earlier graph-by-edges work of Lupanov also belongs in a full historical investigation. The original Lupanov work and the full Wright successor proof were not inspected here. The absence of an asymptotic formula on an OEIS page is not evidence of novelty.
\end{enumerate}

\subsection{Related models and the symmetric-polynomial interpretation}
Greenhill and McKay \cite{greenhill} enumerate vertex-labelled sparse multigraphs with a prescribed degree sequence under a maximum-degree condition. That does not automatically perform the degree-sequence summation and unlabelled reduction needed here. Cameron, Prellberg, and Stark \cite{cameron} use unlabelled vertices but labelled edges in their multigraph interpretation of 2-covers. Dividing such counts by $m!$ without controlling edge symmetries is not justified.

Komiske, Metodiev, and Thaler \cite{komiske} discuss these sequences in the enumeration of symmetric kinematic polynomials. The proposed identification with the relevant polynomial counts, including the OEIS A226919 comparison, is conjectural. Their stable-particle-number limit is different from the present $m\to\infty$ graph-by-edges limit. Theorems in this report concern graph counts and do not assume that conjectural interpretation.

The strongest justified conclusion is therefore a rigorous refined logarithmic enumeration, an unlabelled connectedness ratio, and a bounded-error implicit inverse, together with a separate relative labelled rare-event theorem. A relative equivalent for the unlabelled counts remains a further question.

\begin{thebibliography}{99}
\bibitem{oeisA} The OEIS Foundation, \emph{A307316: Number of unlabeled leafless loopless multigraphs with $n$ edges}. Entry by P. T. Komiske; extended table by A. Howroyd. \url{https://oeis.org/A307316}. Data: \url{https://oeis.org/A307316/b307316.txt}. Checked 2 October 2026.
\bibitem{oeisC} The OEIS Foundation, \emph{A307317: Number of unlabeled connected leafless loopless multigraphs with $n$ edges}. Entry by E. M. Metodiev; extended table by A. Howroyd. \url{https://oeis.org/A307317}. Data: \url{https://oeis.org/A307317/b307317.txt}. Checked 2 October 2026.
\bibitem{wright72} E. M. Wright, The number of unlabelled graphs with many nodes and edges, \emph{Bulletin of the American Mathematical Society} \textbf{78} (1972), 1032--1034, Theorem 5. \url{https://doi.org/10.1090/S0002-9904-1972-13097-0}.
\bibitem{wrighttext} Online reproduction of Wright's 1972 primary text, inspected 2 October 2026. \url{https://paperzz.com/doc/7160301/is-the-number-of-graphs-with-n-unlabelled-nodes-and}. This is a text reproduction, not a verified original scan.
\bibitem{wrightabstract} E. M. Wright, Unlabelled graphs with many nodes and edges, abstract in \emph{Notices of the American Mathematical Society} \textbf{19} (January 1972), A-3. Official issue: \url{https://www.ams.org/journals/notices/197201/197201FullIssue.pdf}. Indexed abstract text checked; original PDF not retrieved.
\bibitem{wright74} E. M. Wright, Graphs on Unlabelled Nodes with a Large Number of Edges, \emph{Proceedings of the London Mathematical Society} (3) \textbf{28} (1974), 577--594. \url{https://doi.org/10.1112/plms/s3-28.4.577}. Metadata checked; full proof not inspected.
\bibitem{gilbert} E. N. Gilbert, Enumeration of labelled graphs, \emph{Canadian Journal of Mathematics} \textbf{8} (1956), 405--411, especially p. 408, equation (9). \url{https://doi.org/10.4153/CJM-1956-046-2}.
\bibitem{greenhill} C. Greenhill and B. D. McKay, Asymptotic enumeration of sparse multigraphs with given degrees. \url{https://arxiv.org/abs/1303.4218}.
\bibitem{cameron} P. J. Cameron, T. Prellberg, and D. Stark, Asymptotic enumeration of 2-covers and line graphs. \url{https://arxiv.org/abs/0707.0664}.
\bibitem{komiske} P. T. Komiske, E. M. Metodiev, and J. Thaler, Cutting Multiparticle Correlators Down to Size. \url{https://arxiv.org/abs/1911.04491}, especially Section V and Table IV.
\end{thebibliography}
\end{document}
